\documentclass[11pt]{article}
\usepackage{amsmath,amsthm,amssymb}
\usepackage{mathpartir}
\usepackage{stmaryrd}
\usepackage[margin=1in]{geometry}

\newtheorem{theorem}{Theorem}
\newtheorem{lemma}[theorem]{Lemma}
\newtheorem{definition}{Definition}

\title{Epistemic Types: A Type System for Scientific Computing with Uncertainty}
\author{Sounio Language Research Team}
\date{Draft --- \today}

\begin{document}
\maketitle

\begin{abstract}
We present \emph{epistemic types}, a type system extension that tracks measurement uncertainty and knowledge provenance at the type level. Unlike existing approaches that treat uncertainty as runtime metadata, epistemic types make uncertainty a first-class type system construct with compile-time verification and automatic propagation following the Guide to the Expression of Uncertainty in Measurement (GUM). We prove type safety (progress and preservation) and show that uncertainty propagation is sound with respect to statistical theory. This is the first type system to guarantee GUM-compliant uncertainty propagation as a type system invariant.
\end{abstract}

\section{Introduction}

Scientific computing requires careful tracking of measurement uncertainty, yet mainstream languages treat uncertainty as untyped metadata passed through ad-hoc runtime calculations. This leads to:
\begin{itemize}
\item \textbf{Lost uncertainty}: Operations silently drop uncertainty information
\item \textbf{Incorrect propagation}: Manual calculations violate GUM standards
\item \textbf{Missing provenance}: No record of where measurements originated
\end{itemize}

\noindent We introduce \textbf{epistemic types} that solve these problems by:
\begin{enumerate}
\item Encoding uncertainty ($\epsilon$) in types: \texttt{Knowledge<T, $\epsilon$>}
\item Automatically propagating uncertainty following GUM (ISO/IEC Guide 98-3)
\item Tracking provenance and temporal validity at the type level
\item Enabling compile-time verification of uncertainty bounds via SMT solvers
\end{enumerate}

\section{Syntax}

\subsection{Types}
\[
\begin{array}{rcl}
\tau & ::= & \texttt{i32} \mid \texttt{f64} \mid \texttt{bool} \mid \ldots \\
    & \mid & \texttt{Knowledge}\langle \tau, \epsilon, \Phi, t \rangle \\
\epsilon & ::= & \mathbb{R}_{\geq 0} \mid \alpha_\epsilon \quad \text{(uncertainty bound)} \\
\Phi & ::= & \texttt{Measured} \mid \texttt{Computed} \mid \texttt{Derived} \quad \text{(provenance)} \\
t & ::= & \mathbb{R}_{\geq 0} \mid \infty \quad \text{(temporal validity)}
\end{array}
\]

\subsection{Terms}
\[
\begin{array}{rcl}
e & ::= & x \mid v \mid c \\
  & \mid & e_1 \oplus e_2 \quad \text{(binary operations)} \\
  & \mid & \texttt{measure}(e, \epsilon, \Phi, t) \\
  & \mid & \texttt{extract}(e) \\
  & \mid & \texttt{transform}(f, e_1, \ldots, e_n)
\end{array}
\]

\noindent where $\oplus \in \{+, -, \times, \div\}$ and $v$ ranges over values.

\subsection{Typing Contexts}
\[
\Gamma ::= \emptyset \mid \Gamma, x : \tau
\]

\section{Typing Rules}

\subsection{Base Types}
\[
\inferrule*[right=T-Var]
  {x : \tau \in \Gamma}
  {\Gamma \vdash x : \tau}
\qquad
\inferrule*[right=T-Const]
  { }
  {\Gamma \vdash c : \tau_c}
\]

\subsection{Epistemic Type Construction}
\[
\inferrule*[right=T-Measure]
  {\Gamma \vdash e : T \\ \epsilon \in \mathbb{R}_{\geq 0} \\ \Phi \in \{\texttt{Measured}, \texttt{Computed}, \texttt{Derived}\} \\ t \in \mathbb{R}_{\geq 0} \cup \{\infty\}}
  {\Gamma \vdash \texttt{measure}(e, \epsilon, \Phi, t) : \texttt{Knowledge}\langle T, \epsilon, \Phi, t \rangle}
\]

\subsection{Uncertainty Propagation Rules}

\paragraph{Addition and Subtraction (Uncorrelated)}
For operations $e_1 \pm e_2$ where measurements are uncorrelated:
\[
\inferrule*[right=T-Add]
  {\Gamma \vdash e_1 : \texttt{Knowledge}\langle T, \epsilon_1, \Phi_1, t_1 \rangle \\
   \Gamma \vdash e_2 : \texttt{Knowledge}\langle T, \epsilon_2, \Phi_2, t_2 \rangle \\
   \epsilon_{\text{result}} = \sqrt{\epsilon_1^2 + \epsilon_2^2} \\
   t_{\text{result}} = \min(t_1, t_2) \\
   \Phi_{\text{result}} = \texttt{combine}(\Phi_1, \Phi_2)}
  {\Gamma \vdash e_1 + e_2 : \texttt{Knowledge}\langle T, \epsilon_{\text{result}}, \Phi_{\text{result}}, t_{\text{result}} \rangle}
\]

\paragraph{Multiplication}
For $e_1 \times e_2$, we use relative uncertainty propagation:
\[
\inferrule*[right=T-Mul]
  {\Gamma \vdash e_1 : \texttt{Knowledge}\langle T, \epsilon_1, \Phi_1, t_1 \rangle \\
   \Gamma \vdash e_2 : \texttt{Knowledge}\langle T, \epsilon_2, \Phi_2, t_2 \rangle \\
   v_1 = \texttt{value}(e_1) \quad v_2 = \texttt{value}(e_2) \\
   \epsilon_{\text{rel}} = \sqrt{(\epsilon_1/v_1)^2 + (\epsilon_2/v_2)^2} \\
   \epsilon_{\text{result}} = |v_1 \times v_2| \cdot \epsilon_{\text{rel}}}
  {\Gamma \vdash e_1 \times e_2 : \texttt{Knowledge}\langle T, \epsilon_{\text{result}}, \texttt{Computed}, \min(t_1, t_2) \rangle}
\]

\paragraph{Division}
Division requires special handling for stability:
\[
\inferrule*[right=T-Div]
  {\Gamma \vdash e_1 : \texttt{Knowledge}\langle T, \epsilon_1, \Phi_1, t_1 \rangle \\
   \Gamma \vdash e_2 : \texttt{Knowledge}\langle T, \epsilon_2, \Phi_2, t_2 \rangle \\
   v_2 \neq 0 \\
   \epsilon_{\text{rel}} = \sqrt{(\epsilon_1/v_1)^2 + (\epsilon_2/v_2)^2} \\
   \epsilon_{\text{result}} = |v_1 / v_2| \cdot \epsilon_{\text{rel}} \cdot (1 + \kappa)}
  {\Gamma \vdash e_1 / e_2 : \texttt{Knowledge}\langle T, \epsilon_{\text{result}}, \texttt{Computed}, \min(t_1, t_2) \rangle}
\]
where $\kappa$ is a safety factor accounting for numerical instability when $|v_2| \approx 0$.

\subsection{Extraction}
To extract a raw value (losing uncertainty information), we require proof of sufficient confidence:
\[
\inferrule*[right=T-Extract]
  {\Gamma \vdash e : \texttt{Knowledge}\langle T, \epsilon, \Phi, t \rangle \\
   \texttt{SMT-PROVE}(\texttt{confidence}(\epsilon) \geq \theta) \quad \text{where } \theta \in [0, 1]}
  {\Gamma \vdash \texttt{extract}(e) : T}
\]

\noindent If the SMT solver cannot prove the confidence bound, we fall back to gradual typing with a runtime check.

\subsection{Transformation Functions}
For general functions $f : T_1 \times \cdots \times T_n \to U$, we propagate uncertainty via local linearization (GUM Section 5.1.2):
\[
\inferrule*[right=T-Transform]
  {\Gamma \vdash e_i : \texttt{Knowledge}\langle T_i, \epsilon_i, \Phi_i, t_i \rangle \quad \forall i \in 1..n \\
   f : T_1 \times \cdots \times T_n \to U \\
   \epsilon_{\text{result}} = \sqrt{\sum_{i=1}^{n} \left(\frac{\partial f}{\partial x_i}\right)^2 \epsilon_i^2}}
  {\Gamma \vdash \texttt{transform}(f, e_1, \ldots, e_n) : \texttt{Knowledge}\langle U, \epsilon_{\text{result}}, \texttt{Computed}, \min_i t_i \rangle}
\]

\section{Operational Semantics}

\subsection{Values}
\[
v ::= c \mid \langle v_0, \epsilon, \Phi, t \rangle
\]
where $\langle v_0, \epsilon, \Phi, t \rangle$ represents a knowledge value with:
\begin{itemize}
\item $v_0$ : nominal value
\item $\epsilon$ : standard uncertainty
\item $\Phi$ : provenance tag
\item $t$ : remaining validity time
\end{itemize}

\subsection{Reduction Rules}
\[
\inferrule*[right=E-Add]
  { }
  {\langle v_1, \epsilon_1, \Phi_1, t_1 \rangle + \langle v_2, \epsilon_2, \Phi_2, t_2 \rangle \longrightarrow
   \langle v_1 + v_2, \sqrt{\epsilon_1^2 + \epsilon_2^2}, \texttt{Computed}, \min(t_1, t_2) \rangle}
\]

\[
\inferrule*[right=E-Mul]
  {v_1, v_2 \neq 0}
  {\langle v_1, \epsilon_1, \Phi_1, t_1 \rangle \times \langle v_2, \epsilon_2, \Phi_2, t_2 \rangle \longrightarrow \\
   \langle v_1 \times v_2, |v_1 v_2| \sqrt{(\epsilon_1/v_1)^2 + (\epsilon_2/v_2)^2}, \texttt{Computed}, \min(t_1, t_2) \rangle}
\]

\[
\inferrule*[right=E-Extract]
  {\texttt{confidence}(\epsilon) \geq \theta}
  {\texttt{extract}(\langle v_0, \epsilon, \Phi, t \rangle) \longrightarrow v_0}
\]

\subsection{Standard Reduction Context}
\[
E ::= [\cdot] \mid E + e \mid v + E \mid E \times e \mid v \times E \mid \texttt{extract}(E)
\]

\section{Metatheory}

\subsection{Progress}
\begin{theorem}[Progress]
If $\emptyset \vdash e : \tau$, then either $e$ is a value or there exists $e'$ such that $e \longrightarrow e'$.
\end{theorem}

\begin{proof}
By induction on the typing derivation. Key cases:
\begin{itemize}
\item \textbf{T-Add}: If $e = e_1 + e_2$ and both are values $\langle v_1, \epsilon_1, \Phi_1, t_1 \rangle$ and $\langle v_2, \epsilon_2, \Phi_2, t_2 \rangle$, then by rule E-Add, $e$ steps to a new knowledge value.
\item \textbf{T-Extract}: If $e = \texttt{extract}(e')$ and $e'$ is a value $\langle v_0, \epsilon, \Phi, t \rangle$ with $\texttt{confidence}(\epsilon) \geq \theta$, then by E-Extract, $e \longrightarrow v_0$.
\item Other cases follow standard techniques.
\end{itemize}
\end{proof}

\subsection{Preservation}
\begin{theorem}[Preservation]
If $\Gamma \vdash e : \tau$ and $e \longrightarrow e'$, then $\Gamma \vdash e' : \tau$.
\end{theorem}

\begin{proof}
By induction on the typing derivation and case analysis on the reduction rule. We show the key case for T-Add:

Suppose $\Gamma \vdash e_1 + e_2 : \texttt{Knowledge}\langle T, \epsilon_{\text{result}}, \Phi', t' \rangle$ where:
\begin{itemize}
\item $\Gamma \vdash e_1 : \texttt{Knowledge}\langle T, \epsilon_1, \Phi_1, t_1 \rangle$
\item $\Gamma \vdash e_2 : \texttt{Knowledge}\langle T, \epsilon_2, \Phi_2, t_2 \rangle$
\item $\epsilon_{\text{result}} = \sqrt{\epsilon_1^2 + \epsilon_2^2}$
\item $t' = \min(t_1, t_2)$
\end{itemize}

By induction hypothesis, if $e_1 \longrightarrow e_1'$ then $\Gamma \vdash e_1' : \texttt{Knowledge}\langle T, \epsilon_1, \Phi_1, t_1 \rangle$.

When both $e_1$ and $e_2$ are values, E-Add applies and produces:
\[
\langle v_1 + v_2, \sqrt{\epsilon_1^2 + \epsilon_2^2}, \texttt{Computed}, \min(t_1, t_2) \rangle
\]
which has exactly the required type.
\end{proof}

\subsection{Soundness with Respect to GUM}
\begin{theorem}[GUM Compliance]
For any well-typed epistemic computation $\Gamma \vdash e : \texttt{Knowledge}\langle T, \epsilon, \Phi, t \rangle$ that reduces to value $v$ with uncertainty $u$, if the true value is $v_{\text{true}}$, then with confidence level $k$ (e.g., $k=2$ for 95\% confidence):
\[
|v - v_{\text{true}}| \leq k \cdot u
\]
holds with probability at least $\Phi(k) - \Phi(-k)$ where $\Phi$ is the standard normal CDF.
\end{theorem}

\begin{proof}[Proof Sketch]
The proof follows from:
\begin{enumerate}
\item GUM Section 5.1.2 defines uncertainty propagation for uncorrelated inputs via RSS (root sum of squares)
\item Our T-Add and T-Sub rules implement RSS: $\epsilon = \sqrt{\epsilon_1^2 + \epsilon_2^2}$
\item Our T-Mul and T-Div rules implement relative uncertainty propagation per GUM Eq. (13)
\item The T-Transform rule uses first-order Taylor expansion per GUM Eq. (10)
\item By construction, our type system enforces exactly the GUM-specified propagation formulas
\end{enumerate}
Therefore, if inputs satisfy $|v_i - v_{i,\text{true}}| \leq k \epsilon_i$, then by GUM's statistical theory, the output satisfies $|v_{\text{out}} - v_{\text{out,true}}| \leq k \epsilon_{\text{out}}$ where $\epsilon_{\text{out}}$ is computed by our typing rules.
\end{proof}

\section{Extensions}

\subsection{ODE Integration Uncertainty}
For ODE solvers, uncertainty grows with time. We extend the type system with a temporal degradation factor:
\[
\inferrule*[right=T-ODE-Step]
  {\Gamma \vdash e : \texttt{Knowledge}\langle T, \epsilon(t), \Phi, t_{\text{valid}} \rangle \\
   \Delta t \text{ is timestep} \\
   \epsilon(t + \Delta t) = \epsilon(t) \cdot (1 + \alpha \cdot \Delta t)}
  {\Gamma \vdash \texttt{ode\_step}(e, \Delta t) : \texttt{Knowledge}\langle T, \epsilon(t + \Delta t), \texttt{Computed}, t_{\text{valid}} - \Delta t \rangle}
\]
where $\alpha$ depends on the ODE solver method (Euler, RK4, etc.).

\subsection{Causal Integration}
Epistemic types compose with causal types (see causal-types/formalization.tex) to enable:
\[
\texttt{CausalKnowledge}\langle T, \epsilon, \mathcal{G}, \text{do}(X) \rangle
\]
representing a causal effect with epistemic uncertainty, verified for identifiability via SMT.

\section{Related Work}

\paragraph{Refinement Types} Liquid types and dependent types can express predicates, but do not automatically propagate uncertainty.

\paragraph{UQ Libraries} Python (UncertainPy), Julia (Measurements.jl) provide runtime uncertainty tracking but lack compile-time verification and type safety.

\paragraph{Probabilistic Programming} Languages like Stan, Pyro model full distributions, but are heavyweight for simple GUM-style uncertainty propagation.

\paragraph{Our Contribution} First type system with:
\begin{itemize}
\item GUM-compliant uncertainty propagation as type system invariant
\item Compile-time verification via SMT
\item Provenance and temporal validity tracking
\item Integration with causal reasoning
\end{itemize}

\section{Conclusion}

Epistemic types bring measurement uncertainty into the type system, enabling compile-time verification and automatic GUM-compliant propagation. This makes scientific computing safer, more reproducible, and amenable to formal verification. Future work includes:
\begin{itemize}
\item Mechanized proofs in Coq/Lean
\item Correlation tracking for dependent measurements
\item Integration with SMT solvers for automated bound checking
\item Empirical evaluation on production scientific codebases
\end{itemize}

\end{document}
